\documentclass[11pt]{article}
\usepackage{amsmath}
\usepackage{amssymb}
\usepackage{natbib}

\title{SAGAN Spectral Modeling Components}
\author{}
\date{\today}

\begin{document}

\maketitle

\section{Spectral Modeling Components}

This document describes the key spectral modeling components implemented in the SAGAN package for fitting astronomical spectra, particularly active galactic nuclei (AGN) and quasar spectra.

\subsection{Line\_MultiGauss}

The \texttt{Line\_MultiGauss} class models emission lines using multiple Gaussian components, consisting of a core component and optional wind components. The model is defined as:

\begin{equation}
F(\lambda) = F_{\rm core}(\lambda) + \sum_{i=0}^{N-1} F_{\rm wind,i}(\lambda)
\end{equation}

where the core component is:

\begin{equation}
F_{\rm core}(\lambda) = A_{\rm c} \exp\left[-\frac{1}{2}\left(\frac{v(\lambda) - v_{\rm c}}{\sigma_{\rm c}}\right)^2\right]
\end{equation}

and each wind component is:

\begin{equation}
F_{\rm wind,i}(\lambda) = A_{\rm c} \cdot a_{\rm w,i} \exp\left[-\frac{1}{2}\left(\frac{v(\lambda) - v_{\rm c} - \Delta v_{\rm w,i}}{\sigma_{\rm w,i}}\right)^2\right]
\end{equation}

Here, $v(\lambda) = c(\lambda - \lambda_0)/\lambda_0$ is the velocity offset from the line center $\lambda_0$, $A_{\rm c}$ is the core amplitude, $v_{\rm c}$ is the core velocity shift, $\sigma_{\rm c}$ is the core velocity dispersion, and $a_{\rm w,i}$, $\Delta v_{\rm w,i}$, and $\sigma_{\rm w,i}$ are the relative amplitude, velocity offset, and velocity dispersion of the $i$-th wind component, respectively. The number of components $N$ is a user-specified parameter, with $N=1$ corresponding to a single Gaussian profile.

\subsection{Line\_MultiGauss\_doublet}

The \texttt{Line\_MultiGauss\_doublet} class extends the multi-Gaussian formalism to model line doublets (e.g., [O~III] $\lambda\lambda$4959, 5007 or [S~II] $\lambda\lambda$6716, 6731). The model is:

\begin{equation}
F(\lambda) = F_0(\lambda) + F_1(\lambda)
\end{equation}

where each line component follows the multi-Gaussian structure:

\begin{equation}
F_j(\lambda) = F_{{\rm core},j}(\lambda) + \sum_{i=0}^{N-1} F_{{\rm wind},j,i}(\lambda)
\end{equation}

with:

\begin{equation}
F_{{\rm core},j}(\lambda) = A_{{\rm c},j} \exp\left[-\frac{1}{2}\left(\frac{v_j(\lambda) - v_{\rm c}}{\sigma_{\rm c}}\right)^2\right]
\end{equation}

\begin{equation}
F_{{\rm wind},j,i}(\lambda) = A_{{\rm c},j} \cdot a_{\rm w,i} \exp\left[-\frac{1}{2}\left(\frac{v_j(\lambda) - v_{\rm c} - \Delta v_{\rm w,i}}{\sigma_{\rm w,i}}\right)^2\right]
\end{equation}

Here, $v_j(\lambda) = c(\lambda - \lambda_j)/\lambda_j$ for $j = 0, 1$ represents the velocity offset from each line center, and $A_{{\rm c},0}$ and $A_{{\rm c},1}$ are the individual core amplitudes for the two lines. Crucially, all velocity components ($v_{\rm c}$, $\Delta v_{\rm w,i}$) and velocity dispersions ($\sigma_{\rm c}$, $\sigma_{\rm w,i}$) are shared between the two lines, ensuring kinematic coupling while allowing independent amplitudes.

\subsection{Line\_GaussHermite}

The \texttt{Line\_GaussHermite} class models a single emission line with a fourth-order Gauss--Hermite expansion, allowing controlled deviations from a pure Gaussian profile. The model is:

\begin{equation}
F(\lambda) = G(\lambda)\left[1 + h_3 H_3\big(w(\lambda)\big) + h_4 H_4\big(w(\lambda)\big)\right]
\end{equation}

with

\begin{equation}
G(\lambda) = A\exp\left[-\frac{1}{2}w(\lambda)^2\right],
\qquad
w(\lambda) = \frac{v(\lambda)-v_0}{\sigma},
\qquad
v(\lambda)=c\frac{\lambda-\lambda_0}{\lambda_0},
\end{equation}

and Hermite basis functions implemented as:

\begin{equation}
H_3(w)=\frac{2w^3-3w}{\sqrt{3}},
\qquad
H_4(w)=\frac{4w^4-12w^2+3}{\sqrt{24}}.
\end{equation}

Here, $A$ is the line amplitude, $v_0$ is the velocity shift, and $\sigma$ is the velocity dispersion. The coefficients $h_3$ and $h_4$ control asymmetry (skewness-like distortion) and peak/wing shape (kurtosis-like distortion), respectively. In the implementation, an optional clipping mode sets negative model values to zero.

\subsection{Line\_GaussHermite\_doublet}

The \texttt{Line\_GaussHermite\_doublet} class extends the Gauss--Hermite formalism to doublets by summing two kinematically tied line components:

\begin{equation}
F(\lambda)=F_0(\lambda)+F_1(\lambda),
\end{equation}

where for $j\in\{0,1\}$,

\begin{equation}
F_j(\lambda)=A_j\exp\left[-\frac{1}{2}w_j(\lambda)^2\right]\left[1+h_3H_3\big(w_j(\lambda)\big)+h_4H_4\big(w_j(\lambda)\big)\right],
\end{equation}

\begin{equation}
w_j(\lambda)=\frac{v_j(\lambda)-v_0}{\sigma},
\qquad
v_j(\lambda)=c\frac{\lambda-\lambda_j}{\lambda_j}.
\end{equation}

The two amplitudes $A_0$ and $A_1$ are independent, while $v_0$, $\sigma$, $h_3$, and $h_4$ are shared between the lines, enforcing common kinematics and common non-Gaussian line-shape distortions across the doublet. As with the single-line model, optional clipping can set negative values of the total profile to zero.

\subsection{Line\_Absorption}

The \texttt{Line\_Absorption} class models an absorption line as a Gaussian optical-depth profile with optional partial covering of the background source. The observed (normalized) flux is given by:

\begin{equation}
F(\lambda) = 1 - C_f + C_f \; e^{-\tau(\lambda)}
\end{equation}

where the optical depth profile is a Gaussian in velocity:

\begin{equation}
	au(\lambda) = \tau_0 \; \exp\left[-\frac{1}{2}\left(\frac{v(\lambda) - v_0}{\sigma}\right)^2\right],
\end{equation}

with $v(\lambda) = c(\lambda - \lambda_0)/\lambda_0$ the velocity offset from the line center $\lambda_0$. The parameters are:

- $\tau_0$: optical depth at line center (unitless), controls the line strength.
- $v_0$: velocity shift of the line center (km~s$^{-1}$).
- $\sigma$: velocity dispersion of the absorbing gas (km~s$^{-1}$).
- $C_f$: covering fraction (0 to 1). $C_f = 1$ corresponds to full covering, giving $F(\lambda)=e^{-\tau(\lambda)}$; $C_f<1$ models partial coverage where a fraction $1-C_f$ of the background remains unabsorbed.

This formulation is commonly used to fit interstellar, circumgalactic, or intrinsic AGN absorption features where the absorber only partially covers the emission source or where saturated lines require a covering-fraction treatment.

\subsection{IronTemplate}

The \texttt{IronTemplate} class models the complex Fe~II and Fe~III emission from AGN using empirical templates. Two template options are available: the I~Zw~1 template from \citet{Boroson1992} or the Mrk~493 template from \citet{Park2022}. The model is:

\begin{equation}
F(\lambda) = A \cdot T_{\rm Fe}\left[\lambda/(1+z)\right] \otimes G(\sigma)
\end{equation}

where $T_{\rm Fe}$ is the normalized iron template, $A$ is the amplitude, $z$ is the redshift, and $G(\sigma)$ is a Gaussian convolution kernel accounting for the velocity dispersion $\sigma$ of the AGN broad-line region. The convolution accounts for the difference between the target dispersion and the intrinsic dispersion of the template ($\sigma_{\rm intr} \approx 384$~km~s$^{-1}$):

\begin{equation}
\sigma_{\rm conv} = \sqrt{\sigma^2 - \sigma_{\rm intr}^2}
\end{equation}

The template is interpolated to the observed wavelength grid after redshifting and convolution. A lower bound of $\sigma \geq 386$~km~s$^{-1}$ is enforced to prevent unphysical deconvolution.

\subsection{WindowedPowerLaw1D}

The \texttt{WindowedPowerLaw1D} class implements a power-law continuum model that is zero outside a specified wavelength range $[\lambda_{\rm min}, \lambda_{\rm max}]$:

\begin{equation}
F(\lambda) = \begin{cases}
A \left(\frac{\lambda}{\lambda_0}\right)^{-\alpha} & \text{if } \lambda_{\rm min} \leq \lambda \leq \lambda_{\rm max} \\
0 & \text{otherwise}
\end{cases}
\end{equation}

where $A$ is the amplitude, $\lambda_0$ is the reference wavelength (typically fixed), and $\alpha$ is the power-law index. This model is useful for fitting localized continuum regions without affecting other spectral windows, making it particularly suitable for modeling AGN continuum components in specific wavelength ranges.

\subsection{BlackBody}

The \texttt{BlackBody} class models thermal emission using the Planck function:

\begin{equation}
F(\lambda) = \begin{cases}
S \cdot B_\lambda(T) / \max[B_\lambda(T)] & \text{if } \lambda_{\rm min} \leq \lambda \leq \lambda_{\rm max} \\
0 & \text{otherwise}
\end{cases}
\end{equation}

where $B_\lambda(T)$ is the Planck function:

\begin{equation}
B_\lambda(T) = \frac{2hc^2}{\lambda^5} \frac{1}{\exp\left(\frac{hc}{\lambda k_B T}\right) - 1}
\end{equation}

Here, $T$ is the temperature in Kelvin, $S$ is a scaling factor, $h$ is Planck's constant, $c$ is the speed of light, $k_B$ is Boltzmann's constant, and the flux is normalized by its maximum value for numerical convenience. The optional wavelength limits $\lambda_{\rm min}$ and $\lambda_{\rm max}$ allow restricting the blackbody contribution to specific spectral regions, useful for modeling hot dust emission or stellar components in composite spectra.

\subsection{Line Spread Function Convolution}

The \texttt{convolve\_lsf} function accounts for instrumental broadening by convolving spectral models with a Gaussian line spread function (LSF). This is essential when fitting spectroscopic data where the observed line profiles are broadened by the instrument's finite spectral resolution.

The LSF is characterized by a resolving power $R = \lambda/\Delta\lambda$ at a reference wavelength $\lambda_c$, where $\Delta\lambda$ is the full width at half maximum (FWHM) of the instrumental profile. The Gaussian kernel has a wavelength-space standard deviation:

\begin{equation}
\sigma_\lambda = \frac{\lambda_c}{R \cdot 2.3548}
\end{equation}

where the factor 2.3548 converts between FWHM and standard deviation (${\rm FWHM} = 2\sqrt{2\ln 2}\,\sigma \approx 2.3548\,\sigma$).

For any spectral model $F_0(\lambda)$, the convolved model is:

\begin{equation}
F(\lambda) = \int_{-\infty}^{\infty} F_0(\lambda') \, G(\lambda - \lambda'; \sigma_\lambda) \, d\lambda'
\end{equation}

where $G(\lambda; \sigma_\lambda)$ is the normalized Gaussian kernel:

\begin{equation}
G(\lambda; \sigma_\lambda) = \frac{1}{\sqrt{2\pi}\sigma_\lambda} \exp\left[-\frac{\lambda^2}{2\sigma_\lambda^2}\right]
\end{equation}

The convolution is implemented using a Gaussian filter on uniformly-spaced wavelength grids. This operation is applied in-place to model instances, and importantly, it also convolves the analytical derivatives (\texttt{fit\_deriv}) when present, ensuring correct behavior with gradient-based fitting algorithms such as Levenberg--Marquardt.

The convolved models retain metadata indicating the LSF parameters ($\lambda_c$, $R$, $\sigma_\lambda$), which can be inspected using the \texttt{find\_convolved\_submodels} function. This is particularly useful when working with composite models containing multiple convolved components.

\section{References}

\begin{thebibliography}{99}

\bibitem[Boroson \& Green(1992)]{Boroson1992}
Boroson, T.~A., \& Green, R.~F.\ 1992, \apjs, 80, 109

\bibitem[Park et al.(2022)]{Park2022}
Park, D., Woo, J.-H., Bennert, V.~N., et al.\ 2022, \apj, 839, 93

\end{thebibliography}

\end{document}
